\subsection{The product factor and its weighted shift}\label{p:sec:operator}

We now turn the fixed rotation into an operator in a finite factor.
The goal is to express invariance as an identity on ordered pairs,
where the sampling result can act. Use the rotation $\sigma$ fixed
at the end of Section~\ref{p:sec:sampling}, and write
$\mathcal K=\ell^2(\Z)$. The operators $U,\pi(g)$ and their generated
algebra $M$ are those defined at the beginning of the appendix.
This is the usual regular representation of a crossed product by a
measure-preserving transformation
\cite[\S1.4 and Proposition~1.4.5]{AnantharamanPopa};
the class of operators $U\pi(f)$ is also considered in
\cite[Section~5]{dykema-schultz}.
For completeness, we verify the properties of $M$ needed here.

\begin{lemma}\label{p:lem:factor}
The algebra $M$ is a factor of type $\mathrm{II}_1$ with separable
predual. Its elements are essentially bounded measurable fields
$H(x)\in\mathcal B(\mathcal K)$ satisfying
\begin{equation}\label{p:eq:covariance}
 H(\sigma x)=S^{-1}H(x)S\quad\text{almost everywhere},
\end{equation}
and its faithful normal tracial state is
\[
 \tau(H)=\int_X H(x)_{0,0}\,d\mu(x).
\]
\end{lemma}

\begin{proof}
Each generator commutes with scalar multiplication by any function
of $x$. It also commutes with the unitary
$(J\zeta)(x)=S\zeta(\sigma x)$. These commutation relations therefore
hold for every $H\in M$.
Lemma~\ref{lem:scalar-multiplier-fields} therefore represents $H$ by an
essentially bounded measurable field $H(x)$, with products and adjoints
computed fiberwise. Commutation with $J$ gives
\eqref{p:eq:covariance}.

Conversely, let $H$ be an essentially bounded measurable field satisfying
\eqref{p:eq:covariance}.  All its integer-iterate covariance identities
hold on one conull set.  Put $h_r(x)=H(x)_{r,0}$ for $r\in\Z$.
Then
\[
 H(x)_{k+r,k}=h_r(\sigma^k x),
\]
so the $r$th matrix diagonal of $H$ is $U^r\pi(h_r)\in M$.
To recover $H$ from these diagonals, use the constant-field unitaries
$D_t e_k=e^{2\pi i kt}e_k$, $t\in\T$.  The strong integral
\[
 H_N=\int_0^1\mathcal F_N(t)D_tHD_t^*\,dt
 =\sum_{|r|<N}\left(1-\frac{|r|}{N}\right)U^r\pi(h_r)
\]
belongs to $M$, where $\mathcal F_N$ is the Fej\'er kernel.  The map $t\mapsto D_tHD_t^*$ is strongly
continuous and uniformly bounded.  The positivity, unit integral and
concentration at zero of $\mathcal F_N$ imply $H_N\to H$ strongly.
The strong closure of $M$ therefore proves the converse field
characterization.

The displayed state is the vector state associated with the constant
field $e_0$, hence is normal. If $H\geq0$ and $\tau(H)=0$, its
$00$ entry vanishes almost everywhere. Covariance then forces every
diagonal entry to vanish almost everywhere. Positivity, or the
inequality $|H_{k,l}|^2\leq H_{k,k}H_{l,l}$, implies $H=0$.
Thus the state is faithful. For $H,H'\in M$, covariance and the
change of variable $z=\sigma^r x$ give
\[
 \int_X H(x)_{0,r}H'(x)_{r,0}\,d\mu(x)
 =\int_X H'(z)_{0,-r}H(z)_{-r,0}\,d\mu(z).
\]
The sums over $r$ are absolutely integrable by the row and column
$\ell^2$ bounds and Cauchy--Schwarz. Summing proves
$\tau(HH')=\tau(H'H)$.

The rotation $\sigma$ is ergodic by Section~\ref{p:sec:scales}.
If $H$ is central, it commutes with
$\pi(x\mapsto e^{2\pi i x_1})$. In each fiber this last operator has
distinct diagonal entries, because $\theta_1$ is irrational; hence
$H(x)$ is diagonal. Commutation with $U$ makes all its diagonal
entries equal. Equation~\eqref{p:eq:covariance} and ergodicity make
their common value constant almost everywhere. The center of $M$
is therefore scalar. The algebra contains the infinite-dimensional
abelian algebra $\pi(L^\infty(X))$, so this finite factor is not a
matrix algebra and is of type $\mathrm{II}_1$. Finally, the predual
of a von Neumann algebra on a separable Hilbert space is a quotient
of the separable trace class, and is separable.
\end{proof}

\subsubsection*{Weights and cocycles}
Using representatives $x_i\in[0,1)$ for the circle coordinates, put
\begin{equation}\label{p:eq:weights}
 f(x)=\exp\!\left(-\sum_{j\geq1}C_j
                    \sum_{i\in J_j}\lfloor x_i+\theta_i\rfloor\right),
 \qquad T=U\pi(f).
\end{equation}
Since $\sum_i\theta_i\leq2\sum_jK_j^{-1}<\infty$, the first
Borel--Cantelli Lemma shows that only finitely many terms in this
sum are nonzero almost everywhere. Thus $0<f\leq1$ almost
everywhere and $T$ is a nonzero bounded operator.
For every $k\in\Z$ define
\[
 F_k(x)=\exp\!\left(-\sum_{j\geq1}C_j
                    \sum_{i\in J_j}\lfloor x_i+k\theta_i\rfloor\right).
\]
For a fixed $k$, all sufficiently late coordinates have
$|k|\theta_i<1$, and the probability that their floor is nonzero is
$|k|\theta_i$. Borel--Cantelli again proves that $F_k$ is positive
and finite almost everywhere. Taking a countable intersection and
all its $\sigma$-translates gives an invariant Borel set $X_0$ of
full measure on which all these formulas hold and
\begin{equation}\label{p:eq:cocycle}
 F_0=1,\qquad F_{k+n}(x)=F_n(\sigma^k x)F_k(x)
 \quad(k,n\in\Z).
\end{equation}
Indeed, the identity follows coordinatewise from
$\lfloor x_i+(k+n)\theta_i\rfloor
 =\lfloor x_i+k\theta_i\rfloor
  +\lfloor\{x_i+k\theta_i\}+n\theta_i\rfloor$.
Extend each $F_k$ to be positive and finite on $X\setminus X_0$;
its values there will never affect an integral.

For $x,y\in X_0$ with $x\leq y$ coordinatewise, set
\begin{equation}\label{p:eq:ratio}
 R_{x,y}(k)=\frac{F_k(y)}{F_k(x)}\leq1,\qquad k\in\Z.
\end{equation}
The inequality holds for negative as well as positive $k$, since
the floor function is nondecreasing. More precisely, each coordinate
$i\in J_s$ with
\begin{equation}\label{p:eq:threshold-crossing}
 x_i<1-\{k\theta_i\}\leq y_i
\end{equation}
contributes a factor $e^{-C_s}$ to this ratio. All other coordinates
contribute either $1$ or another such factor.

\begin{lemma}\label{p:lem:quasinilpotent}
The weight satisfies $\int_X\log f\,d\mu=-\infty$ and $\Delta(T)=0$.
The operator $T$ is nonzero and norm-quasinilpotent:
$\lim_{n\to\infty}\|T^n\|^{1/n}=0$.
\end{lemma}

\begin{proof}
For $n\geq1$, direct multiplication and \eqref{p:eq:cocycle} give
$T^n=U^n\pi(F_n)$. All the floors defining $F_n$ are nonnegative.
For each fixed $q$, the estimate
$\lfloor x_i+n\theta_i\rfloor\geq n\theta_i-1$ on the first
$q$ blocks therefore gives
\[
 \limsup_{n\to\infty}\|T^n\|^{1/n}
 \leq \exp\!\left(-\sum_{s\leq q}C_s
                              \sum_{i\in J_s}\theta_i\right)
 \leq \exp\!\left(-\sum_{s\leq q}\frac{C_s}{K_s}\right).
\]
Since $C_s\geq K_s^2$, the last expression tends to zero as
$q\to\infty$.
Also, Tonelli's Theorem and the same angle bounds give
\[
 -\int_X\log f\,d\mu
 =\sum_s C_s\sum_{i\in J_s}\theta_i
 \ge\sum_s\frac{C_s}{K_s}=\infty.
\]
Since $|T|=\pi(f)$ and the trace on $\pi(L^\infty(X))$ is integration
against $\mu$, the determinant formula gives $\Delta(T)=0$.
Positivity of $f$ almost everywhere ensures $T\ne0$.
\end{proof}

The uniform norm estimate also gives $\sigma(T)=\{0\}$ and Brown
measure $\delta_0$.

\subsubsection*{From orbit identities to triangular laws}
The sampled measures now allow us to use
Lemma~\ref{lem:compact-transfer}. The order supplies its uniform cap,
and the two marginal $L^2$ bounds supply uniform absolute continuity.

\begin{proposition}\label{p:prop:pair-identity}
Suppose $p\in M$ is a projection and $(1-p)Tp=0$. For every finite
$I\subset\N$,
\begin{equation}\label{p:eq:pair-identity}
 (1-p(y))\diag_{k\in\Z}\bigl(R_{x,y}(k)\bigr)p(x)=0
 \quad\text{for }\lambda_I\text{-almost every }(x,y).
\end{equation}
\end{proposition}

\begin{proof}
Fix this $p$ and a Borel version of its projection field. We may
shrink $X_0$ to a common invariant Borel set of full measure on
which covariance, the cocycle identities, and all the following
countably many power identities hold. Invariance of $\ran p(x)$
under $T(x)$ gives invariance under $T(x)^n$, for each $n\geq1$.
Applying $S^{-n}$ and using covariance yields
\begin{equation}\label{p:eq:orbit-inclusion}
 (1-p(\sigma^n x))
 \diag_{k\in\Z}\bigl(F_n(\sigma^k x)\bigr)p(x)=0.
\end{equation}
Division by $F_n(x)>0$ changes the diagonal entries to
$F_k(\sigma^n x)/F_k(x)$ by \eqref{p:eq:cocycle}.
Proposition~\ref{p:prop:sampling} gives ordered orbit-pair measures
converging to $\lambda_I$ with uniformly $L^2$-bounded marginal
densities. Their first marginals are absolutely continuous, so the
countable collection of orbit identities holds almost everywhere
on each approximant. The marginals of $\lambda_I$ are absolutely
continuous. The uniform $L^2$ bounds give uniform absolute continuity
by Cauchy--Schwarz. Orderedness and \eqref{p:eq:ratio} give the common
cap $B=1$ on every ratio, under both the approximants and the limit.
Lemma~\ref{lem:compact-transfer} now proves
\eqref{p:eq:pair-identity}.
\end{proof}

For later use, the common strict-diagonal lemma applies here: the base
is ergodic, the integral trace is faithful, and $f>0$ almost everywhere.
Thus, for any projection $p\in M$ with $0<p<1$ and $(1-p)Tp=0$,
Lemma~\ref{lem:positive-diagonal} gives
\begin{equation}\label{p:eq:strict-diagonal}
 0<r(x):=p(x)_{0,0}<1\quad\text{for almost every }x.
\end{equation}
